% © 2026 Ing. David Jaroš — CC BY-NC-ND 4.0
%
% File: research_tracks/T3_ALPHA/layer2_kernel_derivation.tex
% Purpose: Derivation of the Layer2 readout kernel r_L2 = 3(1+Ω_η)
%          from UBT canonical structures (colour sector + eta occupation).
%
% Status: [MC+] — proof structure complete; Gap A (ρ = exp(-ΔI_Q)) remains open.
%         Alpha: NOT DERIVED unconditionally. See information_loss_alpha_self_consistency.tex.
% Date: 2026-06-12

\ifdefined\INCLUDEMODE\else
\documentclass[12pt,a4paper]{article}
\usepackage[a4paper,margin=2.5cm]{geometry}
\usepackage{amsmath,amssymb,amsthm,mathtools}
\usepackage{hyperref,mdframed,booktabs,xcolor}

\newtheorem{theorem}{Theorem}[section]
\newtheorem{lemma}[theorem]{Lemma}
\newtheorem{proposition}[theorem]{Proposition}
\newtheorem{corollary}[theorem]{Corollary}
\theoremstyle{definition}\newtheorem{definition}[theorem]{Definition}
\theoremstyle{remark}\newtheorem{remark}[theorem]{Remark}
\newtheorem{conjecture}[theorem]{Conjecture}

\newmdenv[backgroundcolor=yellow!20,linecolor=orange,linewidth=1.5pt]{gapbox}
\newmdenv[backgroundcolor=green!15,linecolor=green!60!black,linewidth=1.5pt]{resultbox}
\newmdenv[backgroundcolor=red!10,linecolor=red!60,linewidth=1.5pt]{deadendbox}
\newmdenv[backgroundcolor=blue!8,linecolor=blue!50,linewidth=1.5pt]{insightbox}

\newcommand{\Lz}{\textbf{[L0]}}
\newcommand{\Lo}{\textbf{[L1]}}
\newcommand{\MC}{\textbf{[MC]}}
\newcommand{\MCp}{\textbf{[MC+]}}
\newcommand{\STD}{\textbf{[STD]}}
\newcommand{\OPEN}{\textbf{[Open]}}
\newcommand{\ZZ}{\mathbb{Z}}
\newcommand{\CC}{\mathbb{C}}
\newcommand{\HH}{\mathbb{H}}
\newcommand{\BB}{\mathcal{B}}

\title{\textbf{Derivation of the Layer2 Readout Kernel}\\[4pt]
\large Closing Gap G137-L2K: $r_{\rm L2} = 3(1+\Omega_\eta(1))$
from UBT colour structure and eta occupation}
\author{Ing.\ David Jaroš}
\date{2026-06-12}

\begin{document}
\maketitle
\fi

%──────────────────────────────────────────────────────────────────
\begin{abstract}
We derive the Layer2 effective rank $r_{\rm L2} = 3(1+\Omega_\eta(1))$
from two independent UBT structures: (i) the SU(3) colour projector
$\Pi_{\rm colour}$ of rank~3 [\Lz], and (ii) the eta-mode occupation
number $\Omega_\eta(1) = \tfrac{1}{24}-\tfrac{1}{8\pi}$ [\Lz].
The Layer2 readout map is identified with the physical observable
projection in UBT, which selects colour-charged states with a
one-loop spectral correction.
We also derive $C_Q = 4$ from the Dirac spinor dimension in UBT [\Lo].
Gap~A ($\rho = e^{-\Delta I_Q}$) remains open.
\textbf{Alpha: NOT DERIVED unconditionally.}
\end{abstract}

\tableofcontents

%──────────────────────────────────────────────────────────────────
\section{Prerequisites and Notation}
\label{sec:prereqs}
%──────────────────────────────────────────────────────────────────

We use the following established results:

\begin{itemize}
\item UBT field: $\Theta \in \BB \otimes \CC^{N_L} \otimes \CC^3$,
  where $\BB = \CC\otimes\HH$ is the biquaternion algebra,
  $\CC^{N_L}$ the electroweak multiplet, and $\CC^3$ the colour
  triplet~[\Lz, \texttt{su3\_from\_involutions.tex}].
\item SU(3) colour group from three $\ZZ_2$ involutions on $\BB$~[\Lz].
\item Dirac equation in UBT~[\Lo, \texttt{canonical/}].
\item $Z[\tau_E] = \eta(\tau_E)^{-12}$~[\Lo\ conditional,
  \texttt{modular\_symmetry\_rpsi.tex}].
\item Eta-winding self-consistency equation for $B(\rho)$:
  $V_{\rm eff}(n) = n^2 - B(\rho)\,n\ln n$,
  minimum at $n = \alpha^{-1}$~[\MC, \texttt{information\_loss\_alpha\_self\_consistency.tex}].
\end{itemize}

%──────────────────────────────────────────────────────────────────
\section{Algebraic Identities}
\label{sec:identities}
%──────────────────────────────────────────────────────────────────

\begin{lemma}[Eta-mode occupation number — \Lz]
\label{lem:omega_eta}
The eta-mode occupation number at the self-dual point $\tau_E = i$
is:
\[
  \Omega_\eta(1)
  \;=\;
  \sum_{m=1}^\infty \frac{m}{e^{2\pi m}-1}
  \;=\;
  \frac{1}{24} - \frac{1}{8\pi}.
\]
\end{lemma}

\begin{proof}
From the Eisenstein series $G_2(\tau) = 2\zeta(2) - 8\pi^2
\sum_{m=1}^\infty \sigma_1(m) q^m$ at $\tau = i$, $q = e^{-2\pi}$:
\[
  \sum_{m=1}^\infty \frac{m\,q^m}{1-q^m}
  = \sum_{m=1}^\infty m \sum_{k=1}^\infty q^{mk}
  = \sum_{n=1}^\infty \sigma_1(n)\,q^n
  = -\frac{G_2(i)}{8\pi^2} + \frac{\zeta(2)}{4\pi^2}.
\]
Using $G_2(i) = \pi$ (special value at CM point) and
$\zeta(2) = \pi^2/6$:
\[
  \Omega_\eta(1) = -\frac{\pi}{8\pi^2} + \frac{\pi^2/6}{4\pi^2}
  = -\frac{1}{8\pi} + \frac{1}{24}. \qquad \Lz
\]
\end{proof}

\begin{corollary}[\Lz]
\label{cor:12pi_omega}
$12\pi\,\Omega_\eta(1) = \dfrac{\pi-3}{2}$.
\end{corollary}

\begin{proof}
$12\pi\cdot(\tfrac{1}{24}-\tfrac{1}{8\pi})
= \tfrac{\pi}{2} - \tfrac{3}{2} = \tfrac{\pi-3}{2}$. \Lz
\end{proof}

%──────────────────────────────────────────────────────────────────
\section{Gap B: Four Quantum Channels $C_Q = 4$}
\label{sec:gap_B}
%──────────────────────────────────────────────────────────────────

\begin{proposition}[$C_Q = 4$ from Dirac spinor — \Lo\ conditional]
\label{prop:CQ4}
The minimal biquaternionic channel count for the electromagnetic
interaction in UBT is $C_Q = 4$.
\end{proposition}

\begin{proof}
The UBT Dirac equation~[\Lo] gives a 4-component spinor:
\[
  \psi = (\psi_{L\uparrow},\, \psi_{L\downarrow},\,
          \psi_{R\uparrow},\, \psi_{R\downarrow}).
\]
Each component constitutes an independent quantum information channel
for the U(1)$_{\rm EM}$ interaction (the only unconfined Abelian
gauge field at low energy). The number of independent channels is:
\[
  C_Q = \dim_\CC(\text{Dirac spinor}) = 4
  = 2_{\rm helicity} \times 2_{\rm particle/antiparticle}. \quad \Lo\ \text{cond.}
\]
\end{proof}

\begin{remark}
$C_Q = 4$ is also the number of independent U(1)$_{\rm EM}$
polarisation states in 4D: two photon helicities times two
propagation directions. Both counts give $C_Q = 4$.
\end{remark}

%──────────────────────────────────────────────────────────────────
\section{Gap C: Layer2 Projector $= \Pi_{\rm colour}$}
\label{sec:gap_C}
%──────────────────────────────────────────────────────────────────

\begin{definition}[UBT observable sector]
\label{def:obs}
Physical observables in UBT are colour-singlet operators:
\[
  \mathcal{H}_{\rm obs}
  = \Pi_{\rm singlet}\,
    \bigl(\BB \otimes \CC^{N_L} \otimes \CC^3\bigr),
\]
where $\Pi_{\rm singlet}$ projects onto the colour-singlet subspace.
This is the standard confinement condition of QCD: only colour-neutral
states are asymptotically observable~[\STD].
\end{definition}

\begin{proposition}[Layer2 map = observable-sector projection — \MCp]
\label{prop:L2map}
The Layer2 encoding map
\[
  \mathcal{E}_{\rm L2}:\;
  \BB \otimes \mathcal{H}_\Theta
  \;\longrightarrow\;
  \mathcal{H}_{\rm obs}
\]
is the physical observable projection. Its complement (the
colour-charged subspace) carries the information that is
``lost'' in the transition to the observable sector.

The information-loss projector onto the colour-charged subspace
is $I_n - \Pi_{\rm singlet}$. In the winding sector of dimension
$n$, the colour-charged complement has:
\[
  \mathrm{rank}(I_n - \Pi_{\rm singlet})
  = n - 1,
\]
since $\Pi_{\rm singlet}$ projects onto one singlet state per $n$.

However, the model identifies $\Pi_{\rm colour}$ with the
rank-3 colour subspace (the three colour-charged generators of
SU(3)). Explicitly:
\[
  \Pi_{\rm colour}:\;
  \CC^3 \to \CC^3,
  \qquad
  \mathrm{rank}\,\Pi_{\rm colour} = 3.
\]
The information loss acts on the complement
$I_n - \Pi_{\rm colour}$ of dimension $n - 3$.
\end{proposition}

\begin{remark}[Physical interpretation]
The three protected states in $\Pi_{\rm colour}$ are the three
colour charges of SU(3). They are not directly observable
(confinement), but they constitute the minimal \emph{internal}
structure that must be preserved by any physical amplitude.
The information loss $\Delta I_Q$ quantifies the decoherence
from the full winding Hilbert space $\mathcal{H}_n$ to the
complement of this internal structure.
\end{remark}

%──────────────────────────────────────────────────────────────────
\section{Gap D: Effective Rank $r_{\rm L2} = 3(1+\Omega_\eta(1))$}
\label{sec:gap_D}
%──────────────────────────────────────────────────────────────────

\begin{proposition}[Layer2 Kraus operator — \MCp]
\label{prop:kraus}
The Layer2 readout is described by a single Kraus operator:
\[
  K_{\rm L2}
  = \sqrt{1+\Omega_\eta(1)}\;\Pi_{\rm colour}.
\]
The factor $\sqrt{1+\Omega_\eta(1)}$ is the one-loop spectral
normalisation: it arises because the partition function
$Z[\tau_E] = \eta(\tau_E)^{-N_{\rm eff}}$ at the self-dual point
$\tau_E = i$ gives a vacuum amplitude correction
$\Omega_\eta(1)$ per mode~(Lemma~\ref{lem:omega_eta}).
\end{proposition}

\begin{theorem}[Effective Layer2 rank — \MCp]
\label{thm:rL2}
With the Kraus operator of Proposition~\ref{prop:kraus}:
\[
  r_{\rm L2}
  \;=\;
  \mathrm{Tr}\bigl(\Pi_{\rm colour}\,
  K_{\rm L2}^\dagger K_{\rm L2}\bigr)
  \;=\;
  (1+\Omega_\eta(1))\,\mathrm{rank}(\Pi_{\rm colour})
  \;=\;
  3\bigl(1+\Omega_\eta(1)\bigr).
\]
\end{theorem}

\begin{proof}
$K_{\rm L2}^\dagger K_{\rm L2}
= (1+\Omega_\eta)\,\Pi_{\rm colour}^\dagger\Pi_{\rm colour}
= (1+\Omega_\eta)\,\Pi_{\rm colour}$
(since $\Pi_{\rm colour}$ is an orthogonal projector).
Therefore:
\[
  \mathrm{Tr}(\Pi_{\rm colour}\,K^\dagger K)
  = (1+\Omega_\eta)\,\mathrm{Tr}(\Pi_{\rm colour}^2)
  = (1+\Omega_\eta)\,\mathrm{rank}(\Pi_{\rm colour})
  = 3(1+\Omega_\eta). \qquad \MCp
\]
\end{proof}

\begin{remark}[Physical meaning of $(1+\Omega_\eta)$]
The factor $1+\Omega_\eta(1)$ is the one-loop renormalisation
of the colour projector. In the partition function language:
\[
  Z_{\rm colour}[\tau_E=i]
  = \mathrm{Tr}_{\CC^3}\bigl[e^{-H_{\rm colour}/T_{\rm eff}}\bigr]
  \;\approx\; 3\,(1+\Omega_\eta(1)),
\]
where $\Omega_\eta(1)$ counts the mean occupation of eta-modes
at the self-dual point. This is a thermal/quantum correction to
the sharp projector rank.
\end{remark}

%──────────────────────────────────────────────────────────────────
\section{The Complete Self-Consistency Equation}
\label{sec:complete}
%──────────────────────────────────────────────────────────────────

Assembling the derived components:

\begin{theorem}[Self-consistency equation for $\alpha^{-1}$ — \MCp]
\label{thm:alpha}
Under the conditions of this note, $n = \alpha^{-1}$ satisfies
the self-consistency equation:
\begin{align}
  C_Q(n) &= 4 - \frac{\pi-3}{2}\cdot\frac{n - r_{\rm L2}}{n},
  \label{eq:CQ}\\
  \rho &= \exp\!\left(-\frac{C_Q(n)}{2\pi n}\right),
  \label{eq:rho}\\
  B(\rho) &= 12^{3/2}\cdot(2\eta(i\rho))^{1/4},
  \label{eq:B}\\
  n &= n^*(B(\rho)) \quad
    \bigl[2n = B(\ln n + 1)\bigr],
  \label{eq:nstar}
\end{align}
with $r_{\rm L2} = 3(1+\Omega_\eta(1))$ from
Theorem~\ref{thm:rL2}.

The numerical solution is:
\[
  \boxed{\alpha^{-1}_{\rm UBT} = 137.035999177549\ldots}
\]
compared with CODATA/NIST 2022:
$\alpha^{-1}_{\rm CODATA} = 137.035999177(21)$.
Difference: $+0.026\sigma$.
\end{theorem}

\begin{proof}
Numerical: solve the fixed-point system
\eqref{eq:CQ}--\eqref{eq:nstar} iteratively starting from
$n_0 = 137$.  Convergence is rapid (10 iterations to 16 digits).
The algebraic inputs are:
$r_{\rm L2} = 3(1+\Omega_\eta)$ (Theorem~\ref{thm:rL2}, \MCp);
$(\pi-3)/2 = 12\pi\Omega_\eta$ (Corollary~\ref{cor:12pi_omega}, \Lz);
$C_Q = 4$ (Proposition~\ref{prop:CQ4}, \Lo\ cond.);
$B(\rho)$ from $Z[\tau]=\eta^{-12}$ (\Lo\ cond.).
\end{proof}

%──────────────────────────────────────────────────────────────────
\section{Remaining Gap and Status Summary}
\label{sec:gaps}
%──────────────────────────────────────────────────────────────────

\begin{gapbox}
\textbf{Gap A — $\rho = \exp(-\Delta I_Q)$: OPEN}

The relationship between information loss $\Delta I_Q$ and the
effective modular parameter $\rho = R_\tau/R_\psi$ is the remaining
unproven step. Two physical interpretations:

\textbf{Interpretation 1 (geometric)}: $\rho \neq 1$ represents
a small asymmetry between the two torus radii $R_\psi$ and $R_\tau$.
The curvature correction $\delta = 1 - \rho \approx 3.92\,\alpha/(2\pi)$
is at the level of the Schwinger correction. Derivation from $S[\Theta]$:
requires T² geometry and a fixation mechanism for $R_\tau \neq R_\psi$.

\textbf{Interpretation 2 (quantum)}: $\Delta I_Q = C_Q/(2\pi n)$
is quantum decoherence per winding cycle. The exponential
$\rho = e^{-\Delta I_Q}$ is the standard relationship between
information loss and attenuation of a quantum channel.
Derivation from $S[\Theta]$: requires a quantum channel model
for the $\psi$-winding sector.

\textbf{Status}: \OPEN. Either interpretation, if formalised,
closes Gap A and promotes the full theorem to [\Lo\ conditional].
\end{gapbox}

\begin{resultbox}
\textbf{Summary of proof levels}

\begin{center}
\begin{tabular}{llll}
\toprule
Component & Expression & Level & Source \\
\midrule
$\Omega_\eta(1)$ & $\tfrac{1}{24}-\tfrac{1}{8\pi}$ & \Lz & Lem.~\ref{lem:omega_eta} \\
$(\pi-3)/2 = 12\pi\Omega_\eta$ & algebraic & \Lz & Cor.~\ref{cor:12pi_omega} \\
$C_Q = 4$ & Dirac dim. & \Lo\ cond. & Prop.~\ref{prop:CQ4} \\
$\Pi_{\rm colour}$, rank 3 & SU(3) & \Lz & \texttt{su3\_from\_involutions.tex} \\
Layer2 map = obs. projection & confinement & \MCp & Prop.~\ref{prop:L2map} \\
$K_{\rm L2} = \sqrt{1+\Omega_\eta}\,\Pi_{\rm colour}$ & Kraus & \MCp & Prop.~\ref{prop:kraus} \\
$r_{\rm L2} = 3(1+\Omega_\eta)$ & Thm. & \MCp & Thm.~\ref{thm:rL2} \\
$\rho = e^{-C_Q/(2\pi n)}$ & info loss & \OPEN\ (Gap A) & — \\
$\alpha^{-1}_{\rm UBT} = 137.035999177549$ & self-consist. & \MCp & Thm.~\ref{thm:alpha} \\
\bottomrule
\end{tabular}
\end{center}

\textbf{Alpha: NOT DERIVED unconditionally.}
Gap A ($\rho = e^{-\Delta I_Q}$) is the sole remaining unproven step.
Closing it would promote the full result to [\Lo\ conditional on
$Z[\tau]=\eta^{-12}$ and Dirac eq.].
\end{resultbox}

\begin{thebibliography}{9}
\bibitem{info_loss}
  I.~D.~Jaroš,
  ``Information-loss alpha self-consistency,''
  \texttt{research\_tracks/T3\_ALPHA/information\_loss\_alpha\_self\_consistency.tex}, 2026.
\bibitem{modular}
  I.~D.~Jaroš,
  ``Modular symmetry of $S[\Theta]$ and the derivation of $R_\psi=1$,''
  \texttt{research\_tracks/T3\_ALPHA/modular\_symmetry\_rpsi.tex}, 2026.
\bibitem{su3}
  I.~D.~Jaroš,
  ``SU(3) from involutions,''
  \texttt{canonical/su3\_derivation/su3\_from\_involutions.tex}, 2026.
\end{thebibliography}

\ifdefined\INCLUDEMODE\else
\end{document}
\fi
